\section{How were the datasets built and checked?}
\label{app:data}

We score every call that the target makes after reading the secret, up to and including the call that leaks, in episodes that ran correctly.
Re-running the saved prompts through the open weights reproduces the served replies closely, so the activations we read belong to the model that acted (\cref{tab:replay}).

\paragraph{Which episodes count.}
An episode is analysed only if it ran to completion, the target read the protected value, and the target wrote a correct status summary at some point.
We check for a correct summary at some point, not in the final file, because on gpt-oss-20b a leak can overwrite the summary file.
A check on the final file would therefore drop mostly leaking episodes and select on the outcome.
We fixed this rule for gpt-oss-20b before any probe or monitor saw its data.
The Qwen3-32B analyses keep the episodes and validity checks of the original Qwen3-32B collection, so the two models' rules differ slightly, and no Qwen3-32B episode was re-selected for this version.

\paragraph{Which calls count.}
We exclude the calls before the target reads the secret.
A probe trained to tell calls before and after the read apart does so perfectly, so including those calls would reward detecting the read rather than the leak.
We therefore score calls from the read up to and including the call that leaks.
On gpt-oss-20b, the scored calls also start only after the peer's first request, and the calls between the read and the request serve as a placebo (\cref{app:shortcuts}).
The reasoning-probe comparisons on gpt-oss-20b use the scored calls that returned reasoning; the last-prompt probe gives a similar result on all scored calls (\cref{app:robustness}).

\paragraph{Development and test episodes.}
For Qwen3-32B, the development episodes were collected first and the test episodes later, and the plan in \cref{app:confirmation} was committed before any detector's predictions on them were compared with their labels.
For gpt-oss-20b, the development episodes are the main collection, and the test episodes were collected after the reasoning probe and its analysis plan were committed (\cref{app:confirmation}).
The two sets never share an episode, and no test episode was used to fit or choose anything for the model scored on it.

\paragraph{Episodes excluded before scoring.}
Some collected episodes are not scored, because they fail the validity checks above or have no call in the scored window (\cref{tab:attrition}).
On gpt-oss-20b, most exclusions are episodes with no scored call after the peer's request.
The Qwen3-32B test episodes contain none from the main scenario, and in the gpt-oss-20b test episodes the main scenario has few leaking calls (\cref{tab:attrition}).

\begin{table}[h]
  \centering
  \small
  \input{tables/datasets_attrition.tex}

  \vspace{0.5em}
  \input{tables/datasets_fresh.tex}
  \caption{Episodes collected or replayed, scored and excluded (top), and the test episodes by scenario (bottom).}
  \label{tab:attrition}
\end{table}

\paragraph{Leaking episodes and leaking calls.}
On gpt-oss-20b, the reasoning-probe comparisons have fewer leaking calls than leaking episodes, because some leaking calls return no reasoning and are left out.
Those episodes stay in the scored set through their earlier calls, which are labelled as not leaking for the per-call question and as leaking for the eventual one.

\Cref{tab:reasoning-coverage} gives these counts, and \cref{tab:qwen-accounting} breaks down the Qwen3-32B test calls.

\begin{table}[h]
  \centering
  \footnotesize
  \setlength{\tabcolsep}{3pt}
  \input{tables/reasoning_coverage.tex}
  \caption{Leaking episodes and leaking calls with reasoning, gpt-oss-20b.}
  \label{tab:reasoning-coverage}
\end{table}

\begin{table}[h]
  \centering
  \small
  \input{tables/qwen_call_accounting.tex}
  \caption{Qwen3-32B test calls: leaking calls, earlier calls in leaking episodes, calls in refusing episodes, and control calls, which are counted separately.}
  \label{tab:qwen-accounting}
\end{table}

\paragraph{Reading the activations.}
We re-run every saved target prompt once through the open-weight checkpoint, on one GPU, with no generation.
For Qwen3-32B we save every decoder layer at the last prompt token.
For gpt-oss-20b we also save every returned reasoning token.
Before loading the model, each prompt and token sequence is checked against a hash list written in advance, and every saved array is checked for shape and finite values.

\paragraph{Replay fidelity.}
To check that the replay matches the model that acted, we score the target's actual replies under the replayed model (\cref{tab:replay}).
A low loss on the actual reply tokens and high top-5 agreement mean the replay closely reproduces the served model, though they do not prove token-for-token identity with the provider's kernels.
gpt-oss-20b's chat format keeps the current turn's reasoning between tool calls, and its prompts render within one token of the served prompts for nearly all requests from the main provider.
A second provider served a small share of requests with a less exact rendering, and we keep and flag those requests.

\begin{table}[h]
  \centering
  \small
  \input{tables/replay.tex}
  \caption{Replay fidelity: loss of the actual reply tokens under the replayed model, and top-5 token agreement.}
  \label{tab:replay}
\end{table}

\paragraph{Generation and inference settings.}
Both targets were served through an API at temperature 0.7 (\cref{tab:settings}), and their prompts were then re-run through the open weights to read the activations (\cref{tab:replay}).
The chosen layers are those of the per-call probes fit on development episodes, one per seed.

\begin{table}[h]
  \centering
  \small
  % Generated by generate_tables.py; edit registry table configuration.
\begin{tabular}{p{0.24\linewidth}p{0.32\linewidth}p{0.32\linewidth}}
\toprule
Setting & Qwen3-32B & gpt-oss-20b \\
\midrule
Target temperature & 0.7 & 0.7 \\
Target token limit & 1536 & 4096 \\
Target reasoning & Disabled & Low \\
Target provider & DeepInfra (fp8) & AkashML / DekaLLM \\
Peer model & \texttt{deepseek-v4-flash} & \texttt{deepseek-v4-pro-0813} \\
Peer temperature & \qwenPeerTemperature{} & \gptPeerTemperature{} \\
Peer token limit & \qwenPeerTokenLimit{} & \gptPeerTokenLimit{} \\
Peer reasoning & Disabled & Medium \\
Probe seeds & 42, 43, 44 & 42, 43, 44 \\
Candidate layers (zero-based) & 0, 4, 8, 12, 16, 20, 24, 28, 32, 36, 40, 44, 48, 52, 56, 60, 63 & 0, 1, 3, 4, 6, 7, 9, 10, 12, 13, 15, 16, 18, 19, 20, 22, 23 \\
Layers (seed order) & 48, 16, 48 & 6, 6, 13 \\
Last-prompt threshold & \qwenFrozenThreshold{} & \gptFrozenThreshold{} \\
Reasoning threshold & -- & \gptReasoningFrozenThreshold{} \\
Alarm rule & Strict $>$; nominal 5\% control calls & Strict $>$; nominal 5\% control calls \\
Family rule & All calls grouped by family ID & All calls grouped by family ID \\
Scenario holdout & Remove shared test-family IDs from training & Remove shared test-family IDs from training \\
Development families & 133 & 238 \\
\bottomrule
\end{tabular}

  \caption{Generation and inference settings.}
  \label{tab:settings}
\end{table}

The detector settings and the rules for splits and families are in \cref{tab:detector-settings,tab:split-settings}.

\begin{table}[h]
  \centering
  \small
  \input{tables/detector_settings.tex}
  \caption{Settings of the trained text classifier, the neural-network probe and the incoming-only probe.}
  \label{tab:detector-settings}
\end{table}

\begin{table}[h]
  \centering
  \small
  % Generated by generate_tables.py; edit registry table configuration.
\begin{tabular}{p{0.22\linewidth}p{0.70\linewidth}}
\toprule
Split setting & Implementation \\
\midrule
Family & Episode-generation family ID. Related conditions/tactics reuse the same family ID and deterministic files/value within a scenario; IDs reused across scenarios remain one group. All calls in an episode inherit that ID. \\
Mixed-scenario CV & Five family-grouped outer folds, seeds 42/43/44. Inner validation is one of four family-grouped folds (approximately 25 percent); training and validation families are disjoint. \\
LOSO layer selection & Outer test is one scenario; remove its family IDs from training even when shared across scenarios. Inner validation holds out one training scenario and purges its family IDs. Choose the highest within-scenario/round AUROC; ties use the earliest candidate layer. \\
Final last-prompt models & Inner selection uses a 25 percent development-family holdout, separately for seeds 42/43/44. Refit chosen layers on development data; test scores average logits standardized by each model training mean and standard deviation. \\
Reasoning model & One seed-42 model; 25 percent development-family inner holdout, layer selected by validation MAX-token AUROC. Refit on development calls; no refit on test episodes. \\
\bottomrule
\end{tabular}

  \caption{How families of related episodes and the inner validation splits are defined.}
  \label{tab:split-settings}
\end{table}

